\chapter{Ultraviolet Divergences and the Supersymmetric Cure}

\lectureoverview{Before introducing the formal mathematics of supersymmetry, we should be precise about the problem it is meant to solve. We shall describe ultraviolet divergences in both position space and momentum space, derive the short-distance behavior of scalar and fermion propagators by dimensional analysis, and estimate the leading loop corrections to a scalar mass. A scalar loop and a fermion loop have the same quadratic cutoff dependence but opposite signs. Supersymmetry becomes compelling when we ask what principle could make their couplings, multiplicities, and masses match well enough for that cancellation to be structural rather than accidental.}

\section{Why Return to Divergences?}

The intention was to begin the machinery of supersymmetry immediately. It is better, however, to delay the abstraction for one lecture and rebuild the motivation. The earlier account of divergences and mass renormalization went by quickly, and supersymmetry will seem like an arbitrary mathematical invention unless we first see exactly what needs to be repaired.

Feynman diagrams enter particle physics in two closely related ways. Most familiarly, they organize scattering amplitudes. We specify incoming and outgoing particles and sum the amplitudes represented by all diagrams of the required order. They also organize the construction of an \emph{effective Lagrangian}. Here the word ``effective'' is not a compliment meaning unusually successful. It means that the description is approximate: inaccessible short-distance details have been eliminated or integrated out, and their influence is encoded in the parameters and interactions that remain.

Every such description has a range of validity. We may specify that range by a shortest trustworthy distance, denoted here by \(\delta\), or by a largest trustworthy momentum,
\begin{equation}
\Lambda\sim \delta^{-1}.
\end{equation}
The two cutoffs are not separate assumptions. They are Fourier-dual descriptions of the same ultraviolet limit.

\section{One Propagator, Two Descriptions}

The propagator is the amplitude for creating or inserting a particle at one spacetime point and removing or detecting it at another. Let the two points be \(x^\mu\) and \(y^\mu\), and write their separation as
\begin{equation}
\Delta^\mu=y^\mu-x^\mu.
\end{equation}
We reserve the lower-case symbol \(\delta\) for the cutoff, so that the separation and the regulator cannot be confused.

Translation invariance implies that a free propagator depends on \(x\) and \(y\) only through \(\Delta\). The same object can be Fourier transformed:
\begin{equation}
G(\Delta)
=\int\frac{d^4k}{(2\pi)^4}\,e^{-ik\cdot\Delta}\,\widetilde G(k).
\end{equation}

When the question ``What are the Fourier conjugates to position?'' is put to the room, the answer is momentum. Consequently, the position-space description of a loop as a propagator returning to its starting point is equivalent to the momentum-space description as an integral over every momentum that can circulate around the loop. Momentum-space rules require one to impose momentum conservation at every vertex; position space keeps the present power-counting argument more transparent.

This equivalence identifies the ultraviolet problem before we calculate anything. In position space, a loop probes the coincident-point limit \(\Delta\to0\). In momentum space, it probes the limit \(|k|\to\infty\). Large momentum means short wavelength, so a short-distance singularity and a high-momentum divergence are the same phenomenon in different coordinates.

The lecture treats the simplest spin classes: scalar fields, fermion fields, and, in the background, gauge fields. Scalars have spin zero and are represented by a one-component field \(\phi(x)\). When the lecture was recorded in 2010, the Higgs had not yet been discovered and was described as the principal scalar expected in particle physics. That historical setting explains the playful phrase ``higgly-wiggly''; the field was physically urgent but not yet experimentally established.

For a real scalar field, the propagator may be written schematically as
\begin{equation}
G_\phi(x,y)=\expect{0|T\!\left\{\phi(x)\phi(y)\right\}|0}
=G_\phi(\Delta).
\end{equation}
The time-ordering convention and normalization will not matter for our power counting. The important fact is that the two field insertions describe creation at one point and annihilation at the other.

\begin{classroomqa}
\audiencequestion{Are \(x\) and \(y\) positions in space or in spacetime? \lecturetimestamp{00:09:31}}
\lecturerresponse{They are spacetime points. Each is a four-vector. The propagator therefore depends on their four-dimensional separation and, in a Lorentz-invariant vacuum, on the corresponding invariant interval.}
\end{classroomqa}

A spin-one-half particle is described by a Dirac field \(\psi_i(x)\). The index \(i\) is a spinor index, not a spatial component, and the propagator is therefore a matrix:
\begin{equation}
S_{ij}(x,y)=\expect{0|T\!\left\{\psi_i(x)\bar\psi_j(y)\right\}|0}
=S_{ij}(\Delta).
\end{equation}

A question to the room supplies the component count: a four-dimensional Dirac field has four complex components. A massless Weyl field may be described with two complex components, so left- and right-handed fields can be treated separately. For the present dimensional argument, only the existence of the spinor matrix structure matters.

Because \(S_{ij}\) is a \(4\times4\) matrix, it can be expanded in the basis of Dirac matrices. We shall suppress those indices at first and restore them when they become necessary. This is not a claim that the matrix structure is absent; it is a deliberate simplification while we establish the ultraviolet power law.

\begin{classroomqa}
\audiencequestion{When does the conjugate or daggered field have to appear in a propagator? \lecturetimestamp{00:14:06}}
\lecturerresponse{For a real neutral scalar in this schematic discussion, the same field contains creation and annihilation operators, so no distinct daggered field is required. For a charged field, particle and antiparticle operators must be distinguished. Thus \(\psi\), \(\psi^\dagger\), and covariantly \(\bar\psi=\psi^\dagger\gamma^0\), play different roles for an electron.}
\end{classroomqa}

The physical Higgs particle is neutral, although the underlying Standard Model Higgs field is an electroweak doublet. The lecture is using the real neutral scalar as its local model; that simplification is sufficient for the ultraviolet estimate.\editorialnote{This distinction prevents the schematic real-scalar example from being mistaken for a complete statement about the electroweak Higgs multiplet.}

\section{Natural Units and the Dimension of a Scalar Field}

Set \(\hbar=c=1\). Length and time then have the same units, while energy, mass, and momentum have units of inverse length:
\begin{equation}
[m]=[E]=[p]=L^{-1}.
\end{equation}
The action appears in the phase \(e^{iS}\), so it is dimensionless. A representative scalar action in four spacetime dimensions is
\begin{equation}
S_\phi=\int d^4x\left[
\frac12\partial_\mu\phi\,\partial^\mu\phi
-\frac12m^2\phi^2
-\frac{\lambda}{4!}\phi^4
+\cdots\right].
\end{equation}
The conventional factors of \(1/2\) and \(1/4!\) do not affect dimensional analysis.

Begin with the kinetic term. Since \([d^4x]=L^4\) and \([\partial_\mu]=L^{-1}\), dimensionlessness gives
\begin{align}
1
&=L^4\left(L^{-1}[\phi]\right)^2,\\
[\phi]&=L^{-1}.
\end{align}
The mass term provides a check:
\begin{equation}
L^4[m^2][\phi^2]
=L^4L^{-2}L^{-2}=1.
\end{equation}

The next question to the room is the dimension of the quartic coupling. It is dimensionless in four spacetime dimensions:
\begin{equation}
L^4[\lambda][\phi]^4=L^4[\lambda]L^{-4}=1
\quad\Longrightarrow\quad [\lambda]=L^0.
\end{equation}
This makes the \(\phi^4\) interaction special in four dimensions. The dimensions of \(\phi^6\) and higher interactions follow by the same rule, but their couplings carry negative powers of mass and therefore announce a more restricted effective-theory role.

Thus the field dimensions already tell us which interactions can appear with dimensionless coefficients. The same method will determine both propagators and the Yukawa coupling without evaluating a detailed Feynman integral.

\section{The Scalar Propagator at Short Distance}

Consider first a massless scalar. There is no intrinsic length scale, so the only available scale is the separation \(\Delta\). Translation invariance leaves only \(\Delta\), Lorentz invariance reduces the dependence to its invariant magnitude, and the product of two scalar fields has dimension \(L^{-2}\). Up to a numerical constant,
\begin{equation}
G_0(\Delta)\sim\frac{1}{\Delta^2}.
\end{equation}
Factors such as \(4\pi^2\) cannot be fixed by dimensional reasoning alone.

A mass introduces the Compton length \(m^{-1}\). Its effect is easiest to understand from the dispersion relation
\begin{equation}
E=\sqrt{\bm p^{,2}+m^2},
\qquad
E=|\bm p|\quad(m=0).
\end{equation}
At \(|\bm p|\gg m\), the mass is negligible. Since high momentum probes short distance, the massive propagator must approach the massless \(1/\Delta^2\) behavior as \(\Delta\to0\). At distances much larger than \(m^{-1}\), propagation is suppressed. The lecture represents that fact schematically by
\begin{equation}
G_m(\Delta)\sim \frac{e^{-m|\Delta|}}{\Delta^2}.
\end{equation}
This is deliberately not an exact formula. In four-dimensional Euclidean space the exact free propagator is proportional to \(mK_1(mr)/r\), with \(r=|\Delta|\); it has the same \(1/r^2\) ultraviolet singularity and an exponential large-distance decay accompanied by a power-law prefactor.\editorialnote{The Bessel-function form supplies the standard precise completion of the lecture's schematic exponential and is not required for its power-counting argument.}

The crucial lesson is simple: mass softens long-distance propagation but does not repair the coincident-point singularity. Near enough to a point, every finite mass becomes irrelevant.

\section{The Scalar Tadpole and Fine Tuning}

The quartic interaction \(\lambda\phi^4\) produces a four-line vertex. Join two of its legs to one another and leave two external scalar legs. The resulting tadpole diagram corrects the coefficient of \(\phi^2\), hence the scalar mass squared.

The preserved frame catches the exact transition from dimensional analysis to this loop correction. The dispersion relation remains on the left, the scalar action and field dimension remain partly visible on the right, and the small tadpole has just appeared in the center.

\begin{figure}[H]
\centering
\includegraphics[width=0.72\textwidth]{lecture_03_figure_02.png}
\lectureframe{00:28:30.560}

\medskip
\begin{tikzpicture}[line width=0.8pt,scale=1.0]
  \draw[dashed] (-2.6,0) -- (0,0) -- (2.6,0);
  \fill (0,0) circle (0.055);
  \draw (0,0) .. controls (-0.85,0.7) and (-0.85,1.6) .. (0,1.6)
                .. controls (0.85,1.6) and (0.85,0.7) .. (0,0);
\end{tikzpicture}
\caption{The board at the scalar self-energy transition, followed by a transcript-backed schematic redraw of the tadpole topology. The screenshot securely shows the loop, the relativistic and massless energy formulas, and a partial action; obscured coefficients and labels are not claimed as visible.}
\end{figure}

At the vertex the scalar may be absorbed and re-emitted while the other two legs propagate from the point back to itself. Apart from numerical and combinatorial factors, the correction is
\begin{equation}
\delta m_\phi^2\sim \lambda G_\phi(0).
\end{equation}
The coincident propagator is infinite in the continuum approximation. An effective theory does not claim validity at arbitrarily short distance, so impose \(|\Delta|\geq\delta\). Then
\begin{equation}
G_\phi(0)\longrightarrow G_\phi(\delta)sim\frac{1}{\delta^2},
\qquad
\delta m_\phi^2\sim\frac{\lambda}{\delta^2}sim\lambda\Lambda^2.
\end{equation}
The divergence is called quadratic because it contains two powers of the momentum cutoff. Momentum-space integration gives the same conclusion: excessive contribution comes from arbitrarily large loop momentum. Whether the shortest scale is eventually set by grand unification, quantum gravity, or something not yet imagined does not change the effective-theory diagnosis.

This is a fine-tuning problem. If \(\delta\) is extremely small and \(\lambda\) is not extraordinarily tiny, the radiative correction is much larger than the scalar mass squared observed at low energy. The bare coefficient and the loop correction must then cancel to implausibly many digits:
\begin{equation}
m_{\mathrm{physical}}^2
=m_{\mathrm{bare}}^2+\delta m_\phi^2.
\end{equation}
The massive propagator does not save the situation. At the distances responsible for the divergence, \(e^{-m|\Delta|}\simeq1\). The ultraviolet loop has forgotten the mass long before it reaches the cutoff.

\section{Fermion Dimensions and Their Propagator}

We now repeat the same dimensional analysis for a Dirac fermion. Its free action is
\begin{equation}
S_\psi=\int d^4x\,
\bar\psi\left(i\gamma^\mu\partial_\mu-m\right)\psi.
\end{equation}
The gamma matrices are numerical and dimensionless. The kinetic term contains only one derivative, so
\begin{align}
1
&=L^4[\bar\psi]L^{-1}[\psi],\\
[\psi]=[\bar\psi]&=L^{-3/2}.
\end{align}
Half-integer powers should not be alarming; fermions bring halves into many places.

\begin{classroomqa}
\audiencequestion{Would it not be easier to derive \([\psi]\) from the mass term? \lecturetimestamp{00:36:02}}
\lecturerresponse{Yes. The mass term gives \(L^4L^{-1}[\bar\psi\psi]=1\), hence \([\bar\psi\psi]=L^{-3}\) and \([\psi]=L^{-3/2}\). The kinetic derivation was used first because it remains meaningful in the massless theory and displays why the fermion mass appears as \(m\), not \(m^2\).}
\end{classroomqa}

Two fermion fields have dimension \(L^{-3}\), so the lecture's compact dimensional shorthand for the massless propagator is
\begin{equation}
S(\Delta)\sim\frac{1}{|\Delta|^3}.
\end{equation}
Restoring its matrix structure gives the more precise numerator required by Lorentz covariance:
\begin{equation}
S_{ij}(\Delta)
\propto
\frac{\left(\gamma_\mu\Delta^\mu\right)_{ij}}
{\left(\Delta^2\right)^2}.
\end{equation}
The numerator adds one power of length, and the denominator therefore contains four powers of the interval; the net dimension remains \(L^{-3}\). As with the scalar, a fermion mass suppresses the propagator at large separation but leaves this short-distance scaling unchanged.

\section{Why a Closed Fermion Loop Is Negative}

We can now ask whether a fermionic diagram might cancel the bad bosonic self-energy. The scalar tadpole has a positive sign, up to positive numerical factors such as powers of \(\pi\) and combinatorial coefficients. A closed fermion loop carries the opposite sign.

\begin{classroomqa}
\audiencequestion{Is the minus sign of the fermion loop related to the earlier discussion of one full rotation? \lecturetimestamp{00:40:21}}
\lecturerresponse{It belongs to the same fermionic sign structure, but the useful proof here proceeds through exchange antisymmetry rather than the geometric orientation of a loop. Cut two loops, exchange the identical particles, and reconnect them. The exchange turns two loops into one and contributes a minus sign only for fermions.}
\end{classroomqa}

The argument is worth giving in full. Suppose a one-loop diagram has sign \(\sigma\), either plus or minus. Two disconnected identical loops have sign
\begin{equation}
\sigma^2=+1.
\end{equation}
Cut both loops at the same time. We now have two identical particles propagating between the cuts. Their ordinary two-particle wavefunction is symmetric for bosons and antisymmetric for fermions:
\begin{align}
\Psi_B(x_1,x_2)&=+\Psi_B(x_2,x_1),\\
\Psi_F(x_1,x_2)&=-\Psi_F(x_2,x_1).
\end{align}
The antisymmetry also explains the Pauli principle: if both fermions are put in the same state, interchange changes nothing physically but must reverse the wavefunction, forcing it to vanish.

Reconnect the cut lines after exchanging the two particles. Topologically, the two loops have become one. The unexchanged two-loop amplitude was positive. The exchange leaves it positive for bosons and reverses it for fermions. Therefore
\begin{equation}
\operatorname{sign}(\hbox{closed boson loop})=+1,
\qquad
\operatorname{sign}(\hbox{closed fermion loop})=-1.
\end{equation}
This conclusion is general: each independent closed fermion loop supplies a factor of \(-1\).

\section{The Yukawa Loop}

Return to the scalar mass. A boson cannot turn into a single fermion because fermion number parity would change. The simplest fermionic correction has the scalar split into a fermion-antifermion pair, which propagates around a loop and recombines. The interaction is a Yukawa coupling,
\begin{equation}
\mathcal L_Y=-g\phi\bar\psi\psi.
\end{equation}
For the Higgs, the fermions may be quarks or leptons. There are two interaction points and hence two factors of \(g\).

The scalar tadpole involved one point, whose overall position could be fixed because translation invariance merely supplies the spacetime volume. The Yukawa loop involves two points. We may fix one point, or equivalently fix their average position, but we must integrate over their relative separation \(\Delta\). Each fermion propagator contributes the short-distance factor \(1/|\Delta|^3\), so
\begin{equation}
\delta m_{\phi,F}^2
\sim
-g^2\int_{|\Delta|\geq\delta}d^4\Delta\,
\frac{1}{|\Delta|^6}.
\end{equation}

After the two propagators, the separation integral, and the two couplings have been written, one last factor is requested from the room. The answer is the minus sign. The two propagators close into a fermion loop, so the entire otherwise positive short-distance integral is multiplied by \(-1\).

Dimensional analysis already determines the divergence. Four powers of length come from \(d^4\Delta\), six occur in the denominator, and the only remaining length scale is the cutoff:
\begin{equation}
\int_{|\Delta|\geq\delta}
\frac{d^4\Delta}{|\Delta|^6}
\sim\frac{1}{\delta^2}.
\end{equation}
In four-dimensional Euclidean radial coordinates this can be made explicit:
\begin{equation}
\int_{\delta}^{R}d^4\Delta\,\frac{1}{|\Delta|^6}
=2\pi^2\int_{\delta}^{R}dr\,\frac{r^3}{r^6}
=\pi^2\left(\frac{1}{\delta^2}-\frac{1}{R^2}\right).
\end{equation}
Spinor traces and propagator normalizations alter the coefficient, not the quadratic power of the cutoff.\editorialnote{The radial integral is a standard explicit version of the dimensional estimate made in the lecture. It is written in Euclidean distance, where the positivity and ultraviolet power counting are transparent.}

\begin{classroomqa}
\audiencequestion{Could reversing or choosing the integration limits change the sign? \lecturetimestamp{00:51:06}}
\lecturerresponse{No. This is not a one-dimensional oriented integral whose limits generate the physical sign. It is an integral over four components, or more cleanly over a positive radial measure. Its ultraviolet magnitude is positive. The negative contribution comes solely from the closed fermion loop.}
\end{classroomqa}

\begin{classroomqa}
\audiencequestion{What is the dimension of the Yukawa coupling \(g\)? \lecturetimestamp{00:51:51}}
\lecturerresponse{It is dimensionless in four dimensions. Since \([\phi]=L^{-1}\) and \([\bar\psi\psi]=L^{-3}\), the product \(\phi\bar\psi\psi\) has dimension \(L^{-4}\), exactly canceled by \(d^4x\).}
\end{classroomqa}

The leading fermionic correction is therefore
\begin{equation}
\delta m_{\phi,F}^2
\sim -C_F\frac{g^2}{\delta^2}
=-C_Fg^2\Lambda^2,
\end{equation}
where \(C_F>0\) contains numerical, spin, and multiplicity factors. The scalar loop has the corresponding form
\begin{equation}
\delta m_{\phi,B}^2
\sim+C_B\frac{\lambda}{\delta^2}
=+C_B\lambda\Lambda^2.
\end{equation}
We have not yet solved the fine-tuning problem, but we have found something indispensable: a contribution with the same ultraviolet degree and the opposite sign.

\section{What the Cancellation Requires}

It would be an extraordinary accident if unrelated couplings happened to obey
\begin{equation}
C_B\lambda=C_Fg^2
\end{equation}
to precisely the accuracy required by a large hierarchy of scales. The lecture abbreviates this as \(g^2=\lambda\), suppressing convention-dependent and multiplicity factors. A genuine solution needs a mathematical structure that forces the relation rather than merely permits it.

That is the supersymmetric motive. Supersymmetry relates bosons and fermions and restricts their couplings to combinations that make many otherwise dangerous divergent terms cancel. The cancellation is then repeated across diagrams because it follows from a symmetry, not because parameters have been adjusted one loop at a time.

Mass matching matters as well. If boson and fermion masses are identical, their full propagators can be paired most closely. If the masses differ slightly, the complete finite answers do not cancel exactly, because the propagators differ at distances comparable to their Compton wavelengths. Their common short-distance asymptotics can nevertheless cancel the divergent part:
\begin{equation}
\mathcal I_B(\Delta)-\mathcal I_F(\Delta)
\quad\hbox{may be finite as}\quad \Delta\to0,
\end{equation}
where \(\mathcal I_B\) and \(\mathcal I_F\) denote the complete bosonic and fermionic loop integrands after vertices, numerators, and multiplicities are included. The propagators themselves have different dimensions and are not being subtracted directly. Even when the ultraviolet terms cancel, a finite residual may remain. This is the opening for approximate or broken supersymmetry. Exact equality gives the cleanest cancellation; sufficiently controlled breaking may still remove the ultraviolet catastrophe while leaving observable mass splittings.

We have therefore reached supersymmetry without yet writing its algebra. The chain of reasoning is complete:
\begin{equation}
\begin{gathered}
G_\phi(\Delta)\sim \Delta^{-2}
\quad\Longrightarrow\quad
\delta m_{\phi,B}^2\sim+C_B\lambda\Lambda^2,\\[0.4em]
S_\psi(\Delta)\sim \Delta^{-3}
\quad\Longrightarrow\quad
\delta m_{\phi,F}^2\sim-C_Fg^2\Lambda^2,\\[0.4em]
\text{symmetry-enforced matching}
\quad\Longrightarrow\quad
\text{cancellation of the quadratic divergence}.
\end{gathered}
\end{equation}
The next task is to understand the abstract symmetry principle that enforces those relations and to see concrete examples of how it acts.
